\section{Data and experimental design}\label{sec:data}

\subsection{Point-in-time data}
Returns, market capitalizations and index membership come from WRDS: CRSP daily stock files for the S\&P~500 (total returns including dividends) and Compustat Global for the KOSPI~200, Nikkei~225 and DAX, from January 2000 to 31 December 2025. Each security carries the dates on which it entered and left the index, so the constituent list on any date can be rebuilt. The risk-free rate is the three-month Treasury bill (FRED \texttt{DTB3}) for the United States and the OECD three-month interbank rate for the other markets.

The original study used a fixed list of stocks chosen from today's index with today's market capitalizations. That design has two biases. Survivorship: firms that left the index during the test period are missing \citep{brown1992survivorship,elton1996survivorship}. Look-ahead in the universe: the selection uses information from after the decision date. Here the universe at each rebalancing date is the 50 largest index members \emph{on that date} by market capitalization \emph{on that date}. Stocks that leave the index drop out at the next rebalancing; new entrants are added and receive LLM forecasts from then on. Over the US test period 64 distinct stocks enter the universe.

\subsection{Period and models}
The test period has to start after the training cutoff of every model used and end before the data end. We use only open-weight models whose developers document the training cutoff \todo{cite docs/R4 table in appendix}: gpt-oss-20b and gpt-oss-120b (June 2024) \citep{openai2025gptoss}, Gemma~3 27B (August 2024) \citep{team2025gemma3} and Llama~3.3 70B (December 2023) \citep{meta2024llama33}. Models with undocumented or later cutoffs are excluded, because a forecast for a period the model has read about is not a forecast \citep{glasserman2023assessing,sarkar2024lookahead}. This leaves 1 September 2024 to 31 December 2025: 35 two-week rebalancing periods per market. The first two quarters (September--December 2024) serve only for the first choice of $\tau$; evaluation covers the 26 rebalancing decisions from January to December 2025 (239 trading days), which include the April 2025 tariff sell-off.

This version of the paper reports the pilot market and model: the S\&P~500 with gpt-oss-20b. \pending{KOSPI 200, Nikkei 225, DAX; gpt-oss-120b, Gemma 3 27B, Llama 3.3 70B; prompt B.}

\subsection{Forecast collection}
For each stock and rebalancing date the model receives the stock's name and ticker, its daily returns over the last ten trading days, and the cap-weighted daily returns of the universe over the same days (prompt A, which keeps the information set of the original study). It is asked for the average daily return in percent over the next ten trading days and returns a single number (Appendix~\ref{app:prompt}). Each prompt is sent $N=20$ times at temperature 1.0. gpt-oss-20b is a reasoning model; we use low reasoning effort. Models are served locally with vLLM \citep{kwon2023efficient} on one NVIDIA H100 GPU. The US collection produced 34{,}762 forecasts for 1{,}750 stock-dates; 0.7\% of responses could not be parsed and are dropped (19.9 valid draws per stock-date on average). Collection took 4.5 GPU-hours.

\subsection{Portfolio construction and evaluation}
$\Sig$ is the Ledoit--Wolf shrinkage estimate \citep{ledoit2004well} from the last 126 trading days. $\delta = 2.5$ \citep{he2002intuition}; Section~\ref{sec:robust} reports $\delta \in \{1.5, 3.5\}$ and a trailing 252-day estimate. Weights are set at the end of each period, drift with prices during the next two weeks, and the one-way turnover against the drifted weights is charged 10 basis points on the first holding day (0, 5 and 25 bp in the appendix).

Comparisons that share $\bq$ and differ only in $\Om$ answer test (b). We also report, as reference points, portfolios that do not use LLM views: equal weight, market-cap weight, long-only mean--variance on the 126-day mean, the Black--Litterman prior alone (no views), Black--Litterman with a momentum view ($\bq$ = trailing ten-day mean return, constant $\Om$), Black--Litterman with the 126-day mean as view and its sampling variance as $\Om$, and an equal-weighted top-ten momentum portfolio. Two LLM portfolios without Black--Litterman isolate the role of the update: mean--variance on $\bq$ directly, and equal weight on the ten stocks with the highest $q_i$.

Differences in annualized Sharpe ratio are tested with a paired circular block bootstrap \citep{politis1992circular,ledoit2008robust} on daily excess returns (block length 10 days, 2{,}000 resamples), which resamples the same dates for both portfolios. We report the difference, its 95\% percentile interval and a two-sided bootstrap $p$-value.

For test (a) each stock-date yields one forecast--outcome pair $(q_i, s_i^2, \bar r_i, \hat\sigma_i^2)$, where $\hat\sigma_i^2$ is the 126-day variance of daily returns known at the decision date. The US pilot has 1{,}700 pairs over 34 dates.
